\documentclass[11pt,a4paper]{article}

\usepackage[utf8]{inputenc}
\usepackage[T1]{fontenc}
\usepackage[brazil]{babel}
\usepackage{lmodern}
\renewcommand{\familydefault}{\sfdefault}
\usepackage{microtype}
\usepackage{needspace}
\usepackage{etoolbox}
\usepackage[a4paper, margin=2.4cm, headheight=14pt]{geometry}
\usepackage{xcolor}
\usepackage{listings}
\usepackage{tabularx}
\usepackage{array}
\usepackage{booktabs}
\usepackage{enumitem}
\usepackage{amsmath, amssymb}
\usepackage[explicit]{titlesec}
\usepackage{fancyhdr}
\usepackage{lastpage}
\usepackage[most]{tcolorbox}
\usepackage{tikz}
\usetikzlibrary{arrows.meta, positioning, shapes.geometric, calc, fit, backgrounds}
\usepackage{xurl}
\usepackage[hidelinks]{hyperref}

% ------------------------------------------------------------
% Cores do tema metropolis, as mesmas dos slides
% ------------------------------------------------------------
\definecolor{mDarkTeal}{HTML}{23373B}
\definecolor{mLightBrown}{HTML}{EB811B}
\definecolor{mLightGreen}{HTML}{14B03D}
\definecolor{mBackground}{HTML}{FAFAFA}
\definecolor{termbg}{RGB}{40,44,52}
\definecolor{termfg}{RGB}{220,220,220}
\definecolor{codekw}{RGB}{235,129,27}
\definecolor{codecom}{RGB}{140,150,160}
\definecolor{codestr}{RGB}{152,195,121}

\hypersetup{pdftitle={Concorrência e sincronização}, pdfauthor={Weslei Ferreira Santos}}

\setlength{\parindent}{0pt}
\setlength{\parskip}{0.55em}
\setlist{itemsep=0.2em, topsep=0.3em}
\setlist[itemize]{label=\textcolor{mDarkTeal}{\small$\blacktriangleright$}}

% ------------------------------------------------------------
% Títulos no estilo dos slides
% ------------------------------------------------------------
\newcommand{\barradesecao}[1]{%
  \begin{tcolorbox}[colback=mDarkTeal, colframe=mDarkTeal, sharp corners, boxrule=0pt,
      left=10pt, right=10pt, top=7pt, bottom=7pt, coltext=white, fontupper=\Large\bfseries]
  \thesection\quad #1
  \end{tcolorbox}}
\titleformat{\section}[block]{}{}{0pt}{\barradesecao{#1}}
\titleformat{name=\section, numberless}[block]{}{}{0pt}{%
  \begin{tcolorbox}[colback=mDarkTeal, colframe=mDarkTeal, sharp corners, boxrule=0pt,
      left=10pt, right=10pt, top=7pt, bottom=7pt, coltext=white, fontupper=\Large\bfseries]
  #1
  \end{tcolorbox}}
\titlespacing*{\section}{0pt}{1.2em}{0.6em}

\titleformat{\subsection}[block]{\large\bfseries\color{mDarkTeal}}{}{0pt}%
  {\thesubsection\quad #1}[\vspace{-\parskip}\vspace{3pt}{\color{mLightBrown}\rule{3.2cm}{1.4pt}}]
\titleformat{name=\subsection, numberless}[block]{\large\bfseries\color{mDarkTeal}}{}{0pt}%
  {#1}[\vspace{-\parskip}\vspace{3pt}{\color{mLightBrown}\rule{3.2cm}{1.4pt}}]
\titlespacing*{\subsection}{0pt}{1.1em}{0.4em}
\pretocmd{\subsection}{\needspace{6\baselineskip}}{}{\PackageWarning{apostila}{pretocmd falhou}}
\pretocmd{\section}{\needspace{16\baselineskip}}{}{\PackageWarning{apostila}{pretocmd falhou}}

\titleformat{\subsubsection}[runin]{\bfseries\color{mDarkTeal}}{}{0pt}{#1.}
\titlespacing*{\subsubsection}{0pt}{0.8em}{0.6em}

% ------------------------------------------------------------
% Cabeçalho e rodapé
% ------------------------------------------------------------
\pagestyle{fancy}
\fancyhf{}
\fancyhead[L]{\small\color{mDarkTeal} MATA58 · Sistemas Operacionais}
\fancyhead[R]{\small\color{mDarkTeal} Concorrência e sincronização}
\fancyfoot[R]{\small\color{mDarkTeal}\thepage/\pageref{LastPage}}
\renewcommand{\headrule}{\color{mLightBrown}\hrule height 0.8pt}
\renewcommand{\footrulewidth}{0pt}

% ------------------------------------------------------------
% Blocos no estilo "block=fill" do metropolis
% ------------------------------------------------------------
\tcbset{
  bloco/.style={enhanced, breakable, sharp corners, boxrule=0pt, frame hidden,
    colback=black!4, colbacktitle=black!11, fonttitle=\bfseries,
    left=8pt, right=8pt, top=5pt, bottom=6pt, toptitle=3pt, bottomtitle=3pt,
    before skip=0.9em, after skip=0.9em}
}
\newtcolorbox{definicao}[1][Definição]{bloco, coltitle=mDarkTeal, title={#1}}
\newtcolorbox{alerta}[1][Atenção]{bloco, coltitle=mLightBrown, title={#1}}
\newtcolorbox{exemplo}[1][Exemplo]{bloco, coltitle=mLightGreen!80!black, title={#1}}
\newtcolorbox{pergunta}[1][Pare e pense]{bloco, colback=mLightBrown!7, colbacktitle=mLightBrown!20,
  coltitle=mLightBrown!70!black, title={#1}}

% ------------------------------------------------------------
% Código e terminal, iguais aos slides
% ------------------------------------------------------------
\lstdefinestyle{terminal}{
    backgroundcolor=\color{termbg},
    basicstyle=\ttfamily\small\color{termfg},
    breaklines=true,
    frame=single,
    rulecolor=\color{termbg},
    framerule=3pt,
    xleftmargin=4pt,
    xrightmargin=4pt,
    columns=fullflexible,
    keepspaces=true,
    showstringspaces=false,
    tabsize=4,
    literate={á}{{\'a}}1 {à}{{\`a}}1 {â}{{\^a}}1 {ã}{{\~a}}1 {é}{{\'e}}1 {ê}{{\^e}}1
             {í}{{\'i}}1 {ó}{{\'o}}1 {ô}{{\^o}}1 {õ}{{\~o}}1 {ú}{{\'u}}1 {ç}{{\c c}}1,
    aboveskip=0.8em,
    belowskip=0.8em
}
\lstdefinestyle{codigo}{
    style=terminal,
    language=C,
    basicstyle=\ttfamily\footnotesize\color{termfg},
    keywordstyle=\color{codekw}\bfseries,
    commentstyle=\color{codecom}\itshape,
    stringstyle=\color{codestr},
    morekeywords={pthread_t, sem_t, size_t, atomic_int, atomic_bool},
    moredelim=[is][\color{codekw}\bfseries]{@}{@}
}
\lstdefinestyle{arquivo}{
    style=codigo,
    numbers=left,
    numberstyle=\tiny\color{codecom},
    numbersep=8pt,
    xleftmargin=22pt,
    framexleftmargin=18pt
}

% ------------------------------------------------------------
% Figuras
% ------------------------------------------------------------
\tikzset{
    proc/.style={draw=mDarkTeal, fill=mDarkTeal!8, rounded corners=3pt, thick,
        minimum width=2.1cm, minimum height=0.8cm, align=center, font=\small},
    procf/.style={proc, draw=mDarkTeal!70, fill=white, minimum width=1.2cm, font=\footnotesize},
    destaque/.style={proc, draw=mLightBrown, fill=mLightBrown!12},
    arquivo/.style={draw=black!50, fill=black!5, thick, minimum width=2.2cm,
        minimum height=1.1cm, align=center, font=\small},
    seta/.style={-{Stealth[length=2.2mm]}, thick, draw=mDarkTeal},
    setad/.style={-{Stealth[length=2.2mm]}, very thick, draw=mLightBrown},
    rotulo/.style={font=\scriptsize, fill=white, inner sep=1.5pt},
    celula/.style={draw=mDarkTeal, minimum size=0.7cm, font=\ttfamily\small, fill=white},
    regiao/.style={draw=mDarkTeal, fill=mDarkTeal!6, minimum width=3.6cm, minimum height=0.75cm,
        font=\small, align=center}
}

\lstdefinestyle{java}{style=codigo, language=Java, morekeywords={synchronized}}
\lstdefinestyle{arquivojava}{style=arquivo, language=Java, morekeywords={synchronized}}
\newcolumntype{R}[1]{>{\raggedright\arraybackslash}p{#1}}

\newcommand{\raiz}{../../../..}
\newcommand{\cmd}[1]{\texttt{#1}}

\begin{document}

% ============================================================
% Capa
% ============================================================
\begin{titlepage}
\thispagestyle{empty}
\vspace*{5cm}
{\color{mDarkTeal}\fontsize{28}{34}\selectfont\bfseries Concorrência\\[0.2em] e sincronização\par}
\vspace{0.8em}
{\Large\color{mDarkTeal!80} Condição de corrida, semáforos, produtor-consumidor e monitores\par}
\vspace{1em}
{\color{mLightBrown}\rule{\linewidth}{1.2pt}}
\vspace{1em}

{\large Apostila da aula de 7 de outubro de 2026\par}
\vspace{3em}
{\large Weslei Ferreira Santos\par}
\vspace{0.6em}
MATA58 · Sistemas Operacionais\\
Instituto de Computação / UFBA

\vfill
{\small\color{black!60} Material-base: Prof. Maycon Leone M. Peixoto. Os programas citados estão no repositório da disciplina e podem ser compilados com \cmd{gcc} ou \cmd{make}.\par}
\end{titlepage}

\tableofcontents
\clearpage

% ============================================================
\section{Antes de começar}
% ============================================================

Na aula de 5/10 vimos que cada processo tem a sua própria memória: o pai somou 0 porque as escritas dos filhos ficaram nas cópias deles, e para levar um resultado de um processo a outro foi preciso um pipe. Nesta aula o problema se inverte. As \textbf{threads} de um mesmo processo dividem a memória. Comunicar fica fácil; o difícil passa a ser impedir que uma thread atrapalhe a outra.

\begin{center}
{\large\bfseries\color{mDarkTeal} Se as threads compartilham a memória, como impedir que uma atrapalhe a outra?}
\end{center}

O caminho da aula é este: primeiro vemos o problema acontecer (a \textbf{condição de corrida}); depois damos nome ao trecho perigoso (\textbf{região crítica}) e estudamos três ferramentas para protegê-lo: o algoritmo de \textbf{Peterson}, os \textbf{semáforos} e os \textbf{monitores}. No meio, o problema clássico do \textbf{produtor-consumidor} mostra para que serve cada ferramenta, e uma simulação visual deixa tudo isso visível.

\subsection{O que você deve conseguir ao final}

\begin{itemize}
    \item Explicar por que \cmd{contador++} em duas threads perde incrementos.
    \item Definir região crítica e exclusão mútua.
    \item Explicar Peterson e o que é espera ocupada.
    \item Dizer o que \cmd{mutex}, \cmd{empty} e \cmd{full} controlam no produtor-consumidor e por que a ordem dos \cmd{sem\_wait} importa.
    \item Explicar \cmd{synchronized}, \cmd{wait} e \cmd{notify} em Java, e por que se usa \cmd{while} antes de \cmd{wait}.
    \item Completar o produtor-consumidor com os semáforos certos.
\end{itemize}

\subsection{Como usar os exemplos}

Na raiz do repositório:

\begin{lstlisting}[style=terminal]
$ gcc -Wall -Wextra examples/dia07/corrida.c -o corrida -pthread && ./corrida
$ gcc -Wall -Wextra examples/dia07/corrida_atomic.c -o corrida-atomic -pthread && ./corrida-atomic
$ gcc -Wall -Wextra examples/dia07/corrida_mutex.c -o corrida-mutex -pthread && ./corrida-mutex
$ gcc -Wall -Wextra peterson-code.c -o peterson -pthread && ./peterson
$ gcc -Wall -Wextra prod_cons.c -o prod_cons -pthread && ./prod_cons
$ make java && timeout 5 java -cp build/java-3b ThreadSync3b
$ make java && timeout 5 java -cp build/java-3s ThreadSync3s
$ make run-caixa               # caixa de correio correta em Java
$ make run-simulation          # janela da simulacao Producao de dinheiro
\end{lstlisting}

A opção \cmd{-pthread} prepara a compilação para usar threads. No Linux atual o programa costuma compilar mesmo sem ela, mas em sistemas mais antigos aparece \cmd{undefined reference to pthread\_create}; use sempre. As mailboxes Java rodam para sempre: \cmd{timeout 5} as encerra depois de 5 segundos (ou use \cmd{Ctrl + C}).

% ============================================================
\section{De processos a threads}
% ============================================================

\subsection{O que é uma thread}

\begin{definicao}
Uma \textbf{thread} é um fluxo de execução dentro de um processo. Um processo pode ter várias threads; todas executam o mesmo programa e enxergam a mesma memória, mas cada uma tem o seu próprio ponto de execução, os seus registradores e a sua pilha.
\end{definicao}

\begin{center}
\begin{tikzpicture}[node distance=0pt,
    reg/.style={draw=mDarkTeal, fill=mDarkTeal!8, minimum width=5.4cm, minimum height=0.6cm, font=\small},
    pil/.style={draw=mLightBrown, fill=mLightBrown!15, minimum width=2.6cm, minimum height=0.9cm, font=\small, align=center}]
    \node[reg] (cod) {código};
    \node[reg, below=of cod] (glo) {variáveis globais};
    \node[reg, below=of glo] (heap) {heap (\texttt{malloc})};
    \node[reg, below=of heap] (arq) {arquivos abertos};
    \node[pil, below=0.25cm of arq.south west, anchor=north west] (p1) {pilha\\thread 1};
    \node[pil, below=0.25cm of arq.south east, anchor=north east] (p2) {pilha\\thread 2};
    \node[draw=black!40, dashed, rounded corners, fit=(cod)(p1)(p2), inner sep=0.2cm,
          label={[font=\scriptsize, text=black!60]above:um processo}] {};
    \node[right=0.6cm of glo, font=\footnotesize, align=left] {\textbf{compartilhado}\\por todas as threads};
    \node[right=0.6cm of p2, font=\footnotesize, align=left] {\textbf{de cada thread}};
\end{tikzpicture}
\end{center}

\begin{center}
\begin{tabular}{@{}R{0.45\linewidth}R{0.45\linewidth}@{}}
\toprule
\textbf{Compartilhado pelas threads} & \textbf{Próprio de cada thread} \\
\midrule
Código do programa & Pilha (variáveis locais, chamadas de função) \\
Variáveis globais e \cmd{static} & Registradores e contador de programa \\
Heap (o que foi alocado com \cmd{malloc}) & Estado (executando, pronta, bloqueada) \\
Arquivos e pipes abertos & \\
\bottomrule
\end{tabular}
\end{center}

\subsection{Criando threads em C}

Com a biblioteca POSIX threads (\cmd{pthread}):

\begin{lstlisting}[style=codigo]
#include <pthread.h>

void *somar(void *arg) { ...; return NULL; }   /* o que a thread executa */

pthread_t t1, t2;
pthread_create(&t1, NULL, somar, NULL);        /* cria; a thread ja comeca a rodar */
pthread_create(&t2, NULL, somar, NULL);
pthread_join(t1, NULL);                        /* espera t1 terminar */
pthread_join(t2, NULL);
\end{lstlisting}

\begin{itemize}
    \item \cmd{pthread\_create} recebe o endereço onde guardar o identificador, atributos (\cmd{NULL} = padrão), a função que a thread vai executar e um argumento para ela.
    \item \cmd{pthread\_join} faz quem chamou esperar a thread terminar. É o equivalente do \cmd{waitpid} para threads.
    \item Ao contrário das funções que vimos na aula passada, as funções \cmd{pthread\_*} não devolvem \cmd{-1}: devolvem \cmd{0} em sucesso ou o \textbf{código do erro}.
\end{itemize}

\begin{center}
\begin{tabular}{@{}lll@{}}
\toprule
& \textbf{Processos (5/10)} & \textbf{Threads (7/10)} \\
\midrule
Criar & \cmd{fork()} & \cmd{pthread\_create()} \\
Esperar terminar & \cmd{waitpid()} & \cmd{pthread\_join()} \\
Memória & cópia separada & \textbf{a mesma} \\
Comunicar & pipe & variável compartilhada \\
Risco principal & dados não chegam & dados se \textbf{atropelam} \\
\bottomrule
\end{tabular}
\end{center}

\subsection{Concorrência e paralelismo}

Duas palavras que costumam ser confundidas:

\begin{definicao}
\textbf{Concorrência}: várias tarefas estão em andamento no mesmo período de tempo, e a execução delas se \textbf{intercala}. Basta uma CPU: o sistema alterna entre elas tão rápido que parecem simultâneas.

\textbf{Paralelismo}: várias tarefas executam \textbf{literalmente ao mesmo tempo}, em núcleos diferentes.
\end{definicao}

Todo programa paralelo é concorrente, mas não o contrário. Os problemas desta aula aparecem nos dois casos: numa CPU única, a troca de thread pode acontecer no meio de uma operação; com vários núcleos, as duas threads podem até executar a mesma linha no mesmo instante. O computador usado para testar os exemplos tinha 12 núcleos.

\subsection{Por que usar threads}

\begin{itemize}
    \item \textbf{Responsividade}: um programa com interface gráfica pode continuar respondendo ao usuário enquanto uma thread faz um cálculo demorado. A simulação desta aula faz exatamente isso.
    \item \textbf{Compartilhamento}: as threads já enxergam os mesmos dados; não é preciso pipe nem cópia.
    \item \textbf{Economia}: criar uma thread é bem mais barato do que criar um processo, porque não é preciso montar um novo espaço de endereçamento nem copiar tabelas. Trocar de uma thread para outra do mesmo processo também é mais barato.
    \item \textbf{Aproveitar vários núcleos}: um processo com uma thread só usa um núcleo de cada vez; com várias, pode usar vários.
\end{itemize}

\subsection{Threads do usuário e threads do núcleo}

Uma thread pode ser gerenciada de dois jeitos:

\begin{itemize}
    \item \textbf{Threads do usuário}: uma biblioteca, dentro do processo, cria e alterna as threads. O núcleo do SO só enxerga um processo. É rápido, mas se uma thread fizer uma chamada de sistema bloqueante, todas param, e não há como usar vários núcleos.
    \item \textbf{Threads do núcleo}: o próprio SO conhece cada thread e a escalona como se fosse um processo leve. Uma thread bloqueada não trava as outras, e threads diferentes podem rodar em núcleos diferentes.
\end{itemize}

Os modelos que combinam os dois são chamados \textbf{N:1} (várias threads do usuário sobre uma do núcleo), \textbf{1:1} (cada thread do usuário é uma thread do núcleo) e \textbf{M:N} (várias sobre várias). No Linux, a biblioteca \cmd{pthread} (NPTL) usa o modelo \textbf{1:1}: cada \cmd{pthread\_create} cria uma tarefa no núcleo pela chamada de sistema \cmd{clone()}, a mesma família do \cmd{fork()}, só que pedindo para compartilhar a memória em vez de copiá-la. As threads do Java tradicional também são 1:1; a partir do Java 21 existem também as \emph{threads virtuais}, no modelo M:N.

\subsection{A ordem de execução não é sua}

O escalonador do Linux é \textbf{preemptivo}: ele pode tirar uma thread da CPU a qualquer momento, por exemplo quando acaba a sua fatia de tempo ou quando chega uma interrupção. Não há como o programa saber em qual instrução isso vai acontecer. Por isso, quando duas threads usam o mesmo dado, precisamos raciocinar sobre \textbf{todas} as ordens possíveis, não só sobre a mais provável.

% ============================================================
\section{Condição de corrida}
% ============================================================

\subsection{\texttt{contador++} não é uma operação só}

Para nós, \cmd{contador++} parece uma coisa única. Para o processador, são três passos:

\begin{enumerate}
    \item \textbf{ler} o valor da memória para um registrador;
    \item \textbf{somar 1} no registrador;
    \item \textbf{escrever} o registrador de volta na memória.
\end{enumerate}

Entre um passo e outro, o escalonador pode tirar a thread da CPU e colocar outra. Num computador com vários núcleos, as duas threads podem até executar esses passos \emph{ao mesmo tempo}.

\subsection{O que o processador realmente executa}

Compilando o \cmd{corrida.c} sem otimização (\cmd{gcc -O0 -S}), a linha \cmd{contador = contador + 1} vira três instruções de máquina (x86-64):

\begin{lstlisting}[style=terminal]
movl  contador(%rip), %eax     # 1. le contador da memoria para o registrador eax
addl  $1, %eax                 # 2. soma 1 no registrador
movl  %eax, contador(%rip)     # 3. escreve eax de volta na memoria
\end{lstlisting}

Com otimização (\cmd{-O2}), o gcc gera uma instrução só:

\begin{lstlisting}[style=terminal]
addl  $1, contador(%rip)       # soma 1 direto na memoria
\end{lstlisting}

E mesmo assim os incrementos se perdem. Uma instrução só \textbf{não} é o mesmo que uma operação atômica: por dentro, o processador ainda lê o valor, soma e escreve, e outro núcleo pode ler o valor antigo no meio. Para que essa soma seja indivisível entre núcleos, é preciso o prefixo \cmd{lock} (\cmd{lock addl}), que é justamente o que o compilador gera quando usamos \cmd{atomic\_fetch\_add} (Seção~\ref{sec:hardware}).

\subsection{Uma intercalação que perde um incremento}\label{sec:intercalacao}

Contador começa em 0 e cada thread quer somar 1. O esperado é 2.

\begin{center}
\begin{tabular}{@{}clll@{}}
\toprule
\textbf{Passo} & \textbf{Thread A} & \textbf{Thread B} & \textbf{contador na memória} \\
\midrule
1 & lê 0 & & 0 \\
2 & & lê 0 & 0 \\
3 & calcula 1 & & 0 \\
4 & & calcula 1 & 0 \\
5 & escreve 1 & & 1 \\
6 & & escreve 1 & \textbf{1} \\
\bottomrule
\end{tabular}
\end{center}

As duas threads leram 0 antes de qualquer uma escrever. A escrita de B apagou a de A, e um incremento se perdeu.

\begin{definicao}[Condição de corrida]
Situação em que o resultado de um programa depende da \textbf{ordem} em que as threads foram intercaladas. Como essa ordem muda a cada execução, o resultado também muda, e o erro pode aparecer só de vez em quando.
\end{definicao}

\subsection{Quantas intercalações existem}

Cada thread executa três passos em ordem (ler, somar, escrever), mas os passos de uma podem se encaixar de qualquer jeito entre os da outra. O número de intercalações possíveis é o número de maneiras de escolher, entre 6 posições, as 3 da thread A:

\[
\binom{6}{3} = \frac{6!}{3!\,3!} = 20.
\]

O resultado só sai correto (2) quando uma thread lê o valor \emph{depois} que a outra escreveu, e isso exige que uma delas faça os três passos inteiros antes de a outra começar. Das 20 intercalações, só \textbf{2} dão certo (AAABBB e BBBAAA); as outras 18 perdem um incremento.

Na prática, as trocas de thread não acontecem a cada instrução: uma thread costuma executar milhares de voltas do laço antes de ser interrompida. Por isso nem toda volta perde um incremento. Mas com um milhão de voltas por thread, e com as duas rodando em núcleos diferentes ao mesmo tempo, as janelas ruins aparecem centenas de milhares de vezes.

\subsection{Um exemplo do mundo real: a conta bancária}

Imagine uma conta com R\$ 100 e duas operações de saque de R\$ 80 feitas ao mesmo tempo, cada uma por uma thread:

\begin{lstlisting}[style=codigo]
if (saldo >= valor)          /* 1. confere se ha saldo        */
    saldo = saldo - valor;   /* 2. desconta                   */
\end{lstlisting}

Se as duas threads executam o passo 1 antes de qualquer uma executar o passo 2, as duas veem \cmd{saldo = 100}, as duas aprovam o saque e a conta termina com R\$ $-60$, ou com R\$ 20 se uma escrita apagar a outra. O problema não está em nenhuma linha isolada: está no fato de que ``conferir e agir'' precisa acontecer \textbf{sem interrupção}. Esse padrão, chamado \emph{check-then-act}, é uma das formas mais comuns de condição de corrida.

\subsection{Por que o erro é difícil de achar}

\begin{itemize}
    \item Ele é \textbf{intermitente}: depende da ordem das threads, que muda com a carga da máquina, o número de núcleos e até com a presença de um \cmd{printf} de depuração.
    \item Um teste que passou mil vezes \textbf{não prova} que o programa está certo; prova só que a intercalação ruim não apareceu nessas mil vezes.
    \item Às vezes o erro some quando o programa é executado no depurador, porque o depurador muda os tempos. Esse tipo de erro ganhou o apelido de \emph{heisenbug}.
\end{itemize}

Condições de corrida já causaram acidentes reais. A mais citada é a da máquina de radioterapia Therac-25 (1985 a 1987): uma corrida entre a tarefa que lia o teclado do operador e a que configurava o feixe fez a máquina aplicar doses muito acima do previsto em alguns pacientes.

\subsection{O programa corrida.c}

\lstinputlisting[style=arquivo]{\raiz/examples/dia07/corrida.c}

Algumas execuções reais:

\begin{lstlisting}[style=terminal]
$ ./corrida
Contador: 1006933 (esperado: 2000000)
Perdemos 993067 incrementos por causa da condicao de corrida.
$ ./corrida
Contador: 1349235 (esperado: 2000000)
Perdemos 650765 incrementos por causa da condicao de corrida.
$ ./corrida
Contador: 1022204 (esperado: 2000000)
Perdemos 977796 incrementos por causa da condicao de corrida.
\end{lstlisting}

Quase metade dos incrementos se perde, e o número é diferente a cada vez.

\begin{alerta}[Data race e comportamento indefinido]
Em C, duas threads acessando a mesma variável sem sincronização, com pelo menos uma escrevendo, é uma \textbf{data race}, e o padrão diz que o comportamento é \textbf{indefinido}. Não é correto afirmar que ``o resultado só pode ser um pouco menor''. O \cmd{volatile} do programa não conserta nada: ele só impede o compilador de juntar o milhão de somas numa conta única, para que a corrida apareça. Acrescentar \cmd{sleep} também não conserta: só muda a probabilidade.
\end{alerta}

% ============================================================
\section{Região crítica e exclusão mútua}
% ============================================================

\subsection{Definições}

\begin{definicao}
A \textbf{região crítica} é o trecho do código que acessa um recurso compartilhado (no \cmd{corrida.c}, a linha \cmd{contador = contador + 1}).

A \textbf{exclusão mútua} é a garantia de que no máximo \textbf{uma} thread por vez esteja dentro da região crítica.
\end{definicao}

Se só uma thread por vez pode executar ``ler, somar, escrever'', a intercalação da Seção~\ref{sec:intercalacao} se torna impossível e nenhum incremento se perde.

\subsection{O que uma boa solução precisa garantir}\label{sec:condicoes}

Tanenbaum lista quatro condições:

\begin{enumerate}
    \item Duas threads nunca estão ao mesmo tempo na região crítica.
    \item A solução não depende da velocidade das threads nem do número de CPUs.
    \item Uma thread fora da região crítica não pode impedir outra de entrar.
    \item Nenhuma thread espera para sempre para entrar.
\end{enumerate}

\subsection{Por que uma simples variável de trava não funciona}

Uma ideia ingênua é usar uma variável \cmd{trava}: 0 significa livre, 1 significa ocupado.

\begin{lstlisting}[style=codigo]
while (trava == 1) { }   /* espera ficar livre */
trava = 1;               /* ocupa */
/* regiao critica */
trava = 0;               /* libera */
\end{lstlisting}

O problema é o mesmo do \cmd{contador++}: entre \emph{ver} que a trava está livre e \emph{escrever} 1 nela, a outra thread pode fazer o mesmo teste. As duas veem 0, as duas escrevem 1 e as duas entram. A própria trava sofre condição de corrida. Precisamos de um protocolo mais cuidadoso, ou de ajuda do hardware e do sistema operacional.

\subsection{Desabilitar interrupções}

Numa máquina com \textbf{uma} CPU, a thread só perde o processador por causa de uma interrupção (do relógio, de um dispositivo). Se ela desabilitar as interrupções antes da região crítica e reabilitar depois, ninguém a interrompe. Essa solução tem dois problemas graves:

\begin{itemize}
    \item Só o \textbf{núcleo do SO} pode desabilitar interrupções. Dar esse poder a um programa de usuário permitiria que ele travasse a máquina inteira.
    \item Com \textbf{vários núcleos}, desabilitar as interrupções de um núcleo não impede que outro núcleo execute a mesma região crítica.
\end{itemize}

O próprio núcleo ainda usa essa técnica por instantes muito curtos, mas ela não serve para programas comuns.

\subsection{Alternância estrita}

Outra ideia é uma variável \cmd{vez} que diz de quem é a vez de entrar:

\begin{lstlisting}[style=codigo]
/* thread 0 */                         /* thread 1 */
while (vez != 0) { }                   while (vez != 1) { }
/* regiao critica */                   /* regiao critica */
vez = 1;                               vez = 0;
\end{lstlisting}

Ela garante a exclusão mútua, mas \textbf{viola a condição 3}: as threads são obrigadas a se alternar. Se a thread 0 quer entrar duas vezes seguidas e a thread 1 está ocupada com outra coisa, fora da região crítica, a thread 0 fica esperando por uma thread que não está nem tentando entrar. Peterson resolve exatamente esse defeito, combinando \cmd{vez} com as bandeiras \cmd{interessado}.

\subsection{Ajuda do hardware: instruções atômicas}\label{sec:hardware}

Os processadores modernos oferecem instruções que leem e escrevem uma posição de memória \textbf{numa operação indivisível}, mesmo com vários núcleos:

\begin{center}
\begin{tabular}{@{}lR{0.62\linewidth}@{}}
\toprule
\textbf{Instrução} & \textbf{O que faz, de forma indivisível} \\
\midrule
\emph{test-and-set} (TSL) & Lê o valor antigo e escreve 1. Se o antigo era 0, a trava estava livre e agora é minha. \\
\emph{exchange} (\cmd{xchg}) & Troca o conteúdo de um registrador com o de uma posição de memória. \\
\emph{compare-and-swap} (\cmd{cmpxchg}) & Se a memória ainda tem o valor esperado, troca pelo novo; senão, não mexe e avisa. \\
\emph{fetch-and-add} (\cmd{lock xadd}) & Soma um valor e devolve o antigo. \\
\bottomrule
\end{tabular}
\end{center}

Com \emph{test-and-set}, a variável de trava da seção anterior passa a funcionar, porque ``ver que está livre'' e ``ocupar'' viram um passo só. Em C11, isso é o \cmd{atomic\_flag}:

\begin{lstlisting}[style=codigo]
static atomic_flag ocupado = ATOMIC_FLAG_INIT;

void entrar(void) { while (atomic_flag_test_and_set(&ocupado)) { } }  /* gira */
void sair(void)   { atomic_flag_clear(&ocupado); }
\end{lstlisting}

Uma trava assim, em que quem espera fica girando, se chama \textbf{spinlock}. Ela funciona para qualquer número de threads (Peterson, só para duas), mas ainda é espera ocupada.

\subsection{Consertando o corrida.c}

Há três jeitos de consertar o \cmd{corrida.c}, todos testados no repositório:

\begin{center}
\begin{tabular}{@{}lR{0.5\linewidth}c@{}}
\toprule
\textbf{Conserto} & \textbf{Como funciona} & \textbf{Tempo medido} \\
\midrule
\cmd{atomic\_fetch\_add} & A própria soma vira indivisível (\cmd{lock addl}). Não há região crítica: há uma operação atômica. & 0{,}02 s \\
\cmd{pthread\_mutex\_t} & Região crítica com trava. Quem encontra a trava ocupada dorme. & 0{,}10 s \\
spinlock (\cmd{atomic\_flag}) & Região crítica com trava. Quem encontra a trava ocupada gira. & 0{,}10 s \\
\bottomrule
\end{tabular}
\end{center}

Os três dão sempre 2000000. Os tempos são de uma execução na máquina de teste (12 núcleos), para 2 milhões de incrementos; servem só para comparar a ordem de grandeza. A operação atômica é a mais rápida porque cada soma é uma instrução só; as travas têm o custo extra de entrar e sair. Quando a região crítica é uma única operação simples, a operação atômica é a escolha natural; quando a região crítica tem várias linhas, como inserir num buffer, é preciso uma trava.

\lstinputlisting[style=arquivo, firstline=12]{\raiz/examples/dia07/corrida_atomic.c}

\lstinputlisting[style=arquivo, firstline=11]{\raiz/examples/dia07/corrida_mutex.c}

% ============================================================
\section{O algoritmo de Peterson}
% ============================================================

\subsection{A ideia}

Peterson resolve a exclusão mútua para \textbf{duas} threads usando só leituras e escritas comuns na memória (com a ordem garantida). Cada thread tem uma bandeira \cmd{interessado} e existe uma variável \cmd{vez} compartilhada:

\begin{lstlisting}[style=terminal]
/* entrada */
interessado[eu] = verdadeiro          "quero entrar"
vez = outro                           "mas, se houver empate, você primeiro"
enquanto interessado[outro] e vez == outro:
    esperar

/* região crítica */

/* saída */
interessado[eu] = falso               "saí, não quero mais"
\end{lstlisting}

\subsection{Os dois cenários}

\begin{center}
\begin{tabular}{@{}R{0.36\linewidth}R{0.55\linewidth}@{}}
\toprule
\textbf{Situação} & \textbf{O que acontece} \\
\midrule
O outro não quer entrar & \cmd{interessado[outro]} é falso: a condição do \cmd{enquanto} falha e eu entro direto. Não há alternância forçada. \\
Os dois querem entrar ao mesmo tempo & Os dois escrevem em \cmd{vez}. Só um valor fica: o da thread que escreveu \textbf{por último}. Essa thread vê \cmd{vez == outro} e espera; a outra entra. \\
\bottomrule
\end{tabular}
\end{center}

Quando a thread que entrou sai, ela faz \cmd{interessado = falso}, e a que estava esperando passa.

\subsection{Por que Peterson está correto}

Peterson cumpre as condições da Seção~\ref{sec:condicoes}. Um argumento informal:

\subsubsection{Exclusão mútua}
Suponha, por absurdo, que as duas threads estão na região crítica ao mesmo tempo. Então as duas escreveram \cmd{interessado = verdadeiro} e as duas escreveram em \cmd{vez}. A variável \cmd{vez} guarda um único valor: o da thread que escreveu por último, digamos a thread 1, que escreveu \cmd{vez = 0}. Quando a thread 1 testou o \cmd{enquanto}, viu \cmd{interessado[0]} verdadeiro e \cmd{vez == 0}, ou seja, a condição verdadeira, e teria de esperar. Ela só sairia do laço se a thread 0 mudasse \cmd{interessado[0]} para falso, o que só acontece quando a thread 0 \textbf{sai} da região crítica. Contradição: as duas não podem estar dentro ao mesmo tempo.

\subsubsection{Progresso}
Se só uma thread quer entrar, \cmd{interessado[outro]} é falso e ela entra sem esperar. Uma thread que não está interessada não atrapalha ninguém: é a condição 3, que a alternância estrita violava.

\subsubsection{Espera limitada}
Se as duas querem entrar, a que esperou entra assim que a outra sair, porque a outra, ao tentar entrar de novo, escreve \cmd{vez = outro} e cede a vez. Nenhuma thread pode passar na frente da outra duas vezes seguidas enquanto a outra está esperando.

\subsection{Peterson em C}

\lstinputlisting[style=arquivo]{\raiz/peterson-code.c}

\begin{lstlisting}[style=terminal]
$ ./peterson
Regiao critica: thread0
Regiao critica: thread1#######
Contador: 20000 (esperado: 20000)
\end{lstlisting}

\subsubsection{Por que atomics}
Com variáveis comuns, o compilador e o processador podem \textbf{reordenar} leituras e escritas para ganhar desempenho, desde que uma thread sozinha não perceba. Peterson depende de a escrita em \cmd{interessado} acontecer antes da leitura de \cmd{vez} e do \cmd{interessado} do outro. As funções \cmd{atomic\_store} e \cmd{atomic\_load} do C11 garantem essa ordem.

\subsubsection{O que dá errado sem a ordem garantida}
Processadores modernos têm \emph{buffers de escrita}: uma escrita pode ficar guardada no núcleo por um instante antes de chegar à memória que os outros núcleos enxergam, enquanto a leitura seguinte já é feita. Se isso acontecer em Peterson, a sequência pode ser:

\begin{enumerate}
    \item a thread 0 escreve \cmd{interessado[0] = 1}, mas a escrita fica no buffer do núcleo 0;
    \item a thread 1 escreve \cmd{interessado[1] = 1}, que fica no buffer do núcleo 1;
    \item as duas leem o \cmd{interessado} da outra na memória e veem \cmd{0};
    \item as duas acham que a outra não quer entrar, e as duas entram.
\end{enumerate}

As operações atômicas com ordenação \emph{sequencialmente consistente} (o padrão de \cmd{atomic\_load} e \cmd{atomic\_store}) obrigam o processador a esvaziar o buffer antes de seguir, como se todas as operações de todas as threads acontecessem numa única ordem global. Esse acordo entre a linguagem e o hardware sobre a ordem das leituras e escritas se chama \textbf{modelo de memória}.

\subsection{Espera ocupada}

Enquanto espera, a thread fica girando no \cmd{while}: testa, testa, testa de novo. Isso é \textbf{espera ocupada} (\emph{busy waiting}).

\begin{itemize}
    \item A thread gasta CPU sem fazer trabalho útil.
    \item Com uma única CPU, ela pode até atrasar a thread que está na região crítica e precisa sair.
    \item Peterson também só funciona para duas threads; para mais, existem generalizações mais complicadas.
\end{itemize}

Seria melhor que a thread que não pode entrar \textbf{dormisse} e fosse \textbf{acordada} quando a região ficasse livre. É isso que os semáforos oferecem.

\subsection{Quando girar compensa}

A espera ocupada não é sempre ruim. Colocar uma thread para dormir e acordá-la depois exige uma \textbf{troca de contexto}, que custa alguns microssegundos. Se a região crítica é curtíssima e a outra thread está rodando em outro núcleo, girar alguns ciclos pode sair mais barato do que dormir. Por isso o núcleo do Linux usa spinlocks internamente, e implementações de mutex costumam girar um pouco antes de desistir e dormir.

Girar é ruim quando a espera pode ser longa, ou quando há uma CPU só: a thread que gira ocupa justamente o processador de que a outra precisa para terminar.

\begin{alerta}[Inversão de prioridade]
Com espera ocupada e prioridades, pode acontecer de uma thread de alta prioridade girar esperando uma trava que está com uma thread de baixa prioridade, que nunca recebe a CPU para soltá-la. Um problema desse tipo, com travas e prioridades, fez o robô Mars Pathfinder reiniciar repetidamente em Marte, em 1997, até ser corrigido por uma atualização enviada da Terra.
\end{alerta}

% ============================================================
\section{Semáforos}
% ============================================================

\subsection{Dormir e acordar: o despertar perdido}

Para evitar a espera ocupada, imagine duas operações do SO: \cmd{dormir()}, que suspende a thread, e \cmd{acordar(t)}, que a libera. Um consumidor ingênuo faria:

\begin{lstlisting}[style=codigo]
if (count == 0) dormir();      /* buffer vazio: dorme */
/* retira um item */
\end{lstlisting}

e o produtor, depois de inserir, \cmd{if (count == 1) acordar(consumidor);}. O defeito: o consumidor lê \cmd{count == 0} e, \textbf{antes} de chamar \cmd{dormir()}, perde a CPU. O produtor insere, vê \cmd{count == 1} e chama \cmd{acordar}, mas o consumidor ainda não está dormindo, e o aviso se perde. Quando o consumidor volta, dorme, e nunca mais é acordado. O produtor enche o buffer e também dorme. Os dois dormem para sempre.

A causa é a mesma de sempre: ``testar a condição'' e ``dormir'' precisam ser \textbf{uma operação só}. Em 1965, Edsger Dijkstra propôs uma variável que guarda os avisos não usados e faz o teste e a espera de forma indivisível: o \textbf{semáforo}. Os nomes originais das operações vêm do holandês: \textbf{P} (\emph{proberen}, testar) e \textbf{V} (\emph{verhogen}, incrementar).

\subsection{Definição}

\begin{definicao}
Um \textbf{semáforo} é um contador inteiro que nunca fica negativo, com duas operações atômicas:

\cmd{sem\_wait(\&s)} (também chamada \emph{down} ou P): se \cmd{s > 0}, diminui 1 e segue; se \cmd{s == 0}, a thread \textbf{dorme} até que alguém faça \cmd{sem\_post}.

\cmd{sem\_post(\&s)} (também chamada \emph{up} ou V): soma 1 e, se houver thread dormindo em \cmd{s}, acorda uma.
\end{definicao}

O teste, a alteração do contador e a decisão de dormir são feitos \textbf{atomicamente} pelo sistema: não há a corrida que estragou a variável de trava.

\begin{exemplo}[Analogia: estacionamento com placa de vagas]
O valor do semáforo é o número na placa de vagas. Quem chega e vê 0 espera parado na entrada, sem ficar dando voltas no quarteirão (isso seria a espera ocupada). Quem sai, soma uma vaga na placa e libera o próximo da fila.
\end{exemplo}

\subsection{Como o semáforo funciona por dentro}

Um semáforo é, no núcleo do SO, um contador e uma \textbf{fila} de threads que esperam:

\begin{lstlisting}[style=terminal]
sem_wait(s):                           sem_post(s):
    [início da parte indivisível]          [início da parte indivisível]
    se s.valor > 0:                        se a fila de s não está vazia:
        s.valor = s.valor - 1                  tira uma thread da fila
    senão:                                     e a marca como pronta
        coloca esta thread na fila         senão:
        e a bloqueia (não gasta CPU)           s.valor = s.valor + 1
    [fim da parte indivisível]             [fim da parte indivisível]
\end{lstlisting}

A ``parte indivisível'' é garantida pelo núcleo com as ferramentas da Seção~\ref{sec:hardware}: um spinlock curtíssimo, só durante a atualização do contador e da fila. Ou seja, a espera ocupada não desapareceu: ela ficou restrita a poucas instruções, dentro do SO, em vez de durar a região crítica inteira. A espera longa, essa sim, vira bloqueio.

No Linux, a biblioteca ainda faz uma otimização: se o semáforo está livre, \cmd{sem\_wait} resolve tudo com uma operação atômica, sem nem chamar o núcleo. Só quando é preciso dormir é que entra a chamada de sistema \cmd{futex}.

\subsection{Semáforo como mutex}

Um semáforo que começa em \textbf{1} funciona como uma \textbf{chave}:

\begin{lstlisting}[style=codigo]
sem_init(&mutex, 0, 1);    /* 1: a chave esta disponivel */

sem_wait(&mutex);          /* pega a chave (se outra thread tiver, dorme) */
contador = contador + 1;   /* regiao critica */
sem_post(&mutex);          /* devolve a chave */
\end{lstlisting}

Com isso, o \cmd{corrida.c} daria sempre 2000000. Um semáforo assim, que só vale 0 ou 1, é chamado de \textbf{semáforo binário} ou \textbf{mutex} (de \emph{mutual exclusion}). Um semáforo que conta mais do que 1 é um \textbf{semáforo contador}.

\subsection{Mutex e semáforo binário não são idênticos}

Os dois começam livres e deixam passar uma thread por vez, mas há uma diferença conceitual:

\begin{itemize}
    \item Um \textbf{mutex} tem \textbf{dono}: só a thread que travou pode destravar. Em \cmd{pthread\_mutex\_t} (com o tipo de verificação de erros), destravar um mutex que é de outra thread é um erro.
    \item Um \textbf{semáforo} não tem dono: qualquer thread pode fazer \cmd{sem\_post}. Isso é o que permite usar \cmd{full} e \cmd{empty} como \textbf{sinais} entre threads: o produtor faz \cmd{sem\_post(\&full)} e quem faz \cmd{sem\_wait(\&full)} é outra thread.
\end{itemize}

Regra prática: para proteger uma região crítica, use um mutex; para contar recursos ou sinalizar eventos entre threads, use um semáforo. O \cmd{prod\_cons.c} usa um semáforo binário como mutex por fidelidade ao exemplo clássico; com \cmd{pthread\_mutex\_t} funcionaria igual.

\subsection{Funções POSIX}

\begin{center}
\begin{tabular}{@{}lR{0.62\linewidth}@{}}
\toprule
\textbf{Função} & \textbf{O que faz} \\
\midrule
\cmd{sem\_init(\&s, 0, v)} & Cria o semáforo com valor inicial \cmd{v}. O \cmd{0} diz que ele é compartilhado entre threads do mesmo processo. \\
\cmd{sem\_wait(\&s)} & Espera e decrementa. Pode ser interrompida por sinal (\cmd{EINTR}); por isso o \cmd{prod\_cons.c} tem o \emph{helper} \cmd{wait\_sem}, que tenta de novo. \\
\cmd{sem\_post(\&s)} & Incrementa e acorda uma thread que espera. \\
\cmd{sem\_destroy(\&s)} & Libera o semáforo quando ninguém mais o usa. \\
\bottomrule
\end{tabular}
\end{center}

% ============================================================
\section{O problema do produtor-consumidor}
% ============================================================

\subsection{O problema}

Uma thread \textbf{produz} itens e os coloca num buffer de tamanho \cmd{N}; outra \textbf{consome} itens desse buffer. As duas trabalham em ritmos diferentes. Precisamos garantir três coisas:

\begin{enumerate}
    \item As duas nunca mexem no buffer ao mesmo tempo (\textbf{exclusão mútua}).
    \item O produtor não insere com o buffer cheio (\textbf{espera por vaga}).
    \item O consumidor não retira com o buffer vazio (\textbf{espera por item}).
\end{enumerate}

\begin{center}
\begin{tikzpicture}[
    cel/.style={draw=mDarkTeal, minimum width=1.05cm, minimum height=0.8cm, font=\ttfamily\small, fill=white},
    chei/.style={cel, fill=mLightBrown!25}]
    \node[proc] (prod) at (-1.2,0) {Produtor};
    \node[chei] (c0) at (2.0,0) {7};
    \node[chei] (c1) at (3.05,0) {42};
    \node[chei] (c2) at (4.1,0) {13};
    \node[cel] (c3) at (5.15,0) {};
    \node[cel] (c4) at (6.2,0) {};
    \foreach \i in {0,...,4} \node[font=\tiny, below=1pt of c\i] {\i};
    \node[proc] (cons) at (9.4,0) {Consumidor};
    \draw[setad] (prod) -- (c0.west);
    \draw[setad] (c4.east) -- (cons);
    \node[font=\scriptsize\bfseries, text=mLightBrown, above=0.15cm of c3] {\texttt{hi}};
    \node[font=\scriptsize\bfseries, text=mDarkTeal, above=0.15cm of c0] {\texttt{lo}};
\end{tikzpicture}
\end{center}

O buffer é \textbf{circular}: \cmd{hi} é a próxima posição de escrita e \cmd{lo} a próxima de leitura. Os dois índices avançam com \cmd{(i + 1) \% N}, e depois da última posição voltam para a 0. Assim as posições são reaproveitadas e a ordem de chegada é mantida (FIFO).

\subsection{Os três semáforos}

\begin{center}
\begin{tabular}{@{}llR{0.55\linewidth}@{}}
\toprule
\textbf{Semáforo} & \textbf{Começa em} & \textbf{O que controla} \\
\midrule
\cmd{mutex} & 1 & \textbf{Exclusão mútua}: a chave do buffer. \\
\cmd{empty} & \cmd{N} & \textbf{Disponibilidade}: quantas vagas livres existem. \\
\cmd{full} & 0 & \textbf{Disponibilidade}: quantos itens prontos existem. \\
\bottomrule
\end{tabular}
\end{center}

\begin{definicao}[Dois problemas diferentes]
O \cmd{mutex} responde ``\textbf{posso mexer no buffer agora?}''. Os semáforos \cmd{empty} e \cmd{full} respondem ``\textbf{tem o que eu preciso?}''. Confundir os dois é o erro mais comum neste problema.
\end{definicao}

\subsection{Os passos de cada thread}

\begin{center}
\begin{tabular}{@{}cR{0.42\linewidth}R{0.42\linewidth}@{}}
\toprule
\textbf{Passo} & \textbf{Produtor} & \textbf{Consumidor} \\
\midrule
1 & \cmd{sem\_wait(\&empty)}: reserva uma vaga & \cmd{sem\_wait(\&full)}: reserva um item \\
2 & \cmd{sem\_wait(\&mutex)}: pega a chave & \cmd{sem\_wait(\&mutex)}: pega a chave \\
3 & insere no buffer & retira do buffer \\
4 & \cmd{sem\_post(\&mutex)}: devolve a chave & \cmd{sem\_post(\&mutex)}: devolve a chave \\
5 & \cmd{sem\_post(\&full)}: avisa que há mais um item & \cmd{sem\_post(\&empty)}: avisa que há mais uma vaga \\
\bottomrule
\end{tabular}
\end{center}

\begin{exemplo}[Regra para lembrar]
Cada thread \textbf{espera o que consome} e \textbf{avisa o que cria}. O produtor consome vagas e cria itens; o consumidor consome itens e cria vagas.
\end{exemplo}

\subsection{Simulação com N = 3}

Valores logo depois de cada operação terminar:

\begin{center}
\begin{tabular}{@{}lccc@{}}
\toprule
\textbf{Evento} & \textbf{empty} & \textbf{full} & \textbf{count} \\
\midrule
início & 3 & 0 & 0 \\
produz A & 2 & 1 & 1 \\
produz B & 1 & 2 & 2 \\
consome A & 2 & 1 & 1 \\
produz C & 1 & 2 & 2 \\
produz D & 0 & 3 & 3 \\
\bottomrule
\end{tabular}
\end{center}

Se o produtor tentar um quinto item, \cmd{empty} vale 0 e ele dorme em \cmd{sem\_wait(\&empty)}, \textbf{sem segurar o mutex}, até o consumidor liberar uma vaga.

Repare que \cmd{empty + full = N} em todas as linhas. Isso vale sempre que nenhuma operação está em andamento; no meio de uma operação (depois do \cmd{sem\_wait} e antes do \cmd{sem\_post}), a soma pode ficar menor por um instante.

\subsection{O programa prod\_cons.c}

\lstinputlisting[style=arquivo]{\raiz/prod_cons.c}

\begin{lstlisting}[style=terminal]
$ ./prod_cons
Produzido: 29 | ocupacao: 1/20
Produzido: 51 | ocupacao: 2/20
Produzido: 63 | ocupacao: 3/20
...
Fim: 40 itens produzidos e consumidos; buffer vazio.
\end{lstlisting}

A ordem das linhas ``Produzido'' e ``Consumido'' muda a cada execução, mas os itens sempre são consumidos na mesma ordem em que foram produzidos.

\subsection{Por que a ordem dos sem\_wait importa}

Suponha que o consumidor pegasse a chave \emph{antes} de esperar por item:

\begin{lstlisting}[style=codigo]
wait_sem(&mutex);   /* pega a chave...     */
wait_sem(&full);    /* ...e dorme com ela  */
\end{lstlisting}

Com o buffer vazio:

\begin{enumerate}
    \item O consumidor pega o \cmd{mutex}.
    \item Dorme em \cmd{full}, \textbf{segurando a chave}.
    \item O produtor tenta pegar o \cmd{mutex} para inserir e também dorme.
\end{enumerate}

\begin{center}
\begin{tikzpicture}[node distance=3cm]
    \node[proc] (c) {Consumidor\\\scriptsize tem o \texttt{mutex}};
    \node[proc, right=of c] (p) {Produtor\\\scriptsize quer o \texttt{mutex}};
    \draw[setad, bend left=22] (c) to node[rotulo, above] {espera um item do produtor} (p);
    \draw[setad, bend left=22] (p) to node[rotulo, below] {espera a chave do consumidor} (c);
\end{tikzpicture}
\end{center}

\begin{alerta}[Deadlock]
Cada thread espera algo que só a outra pode fornecer. Nenhuma avança, e o programa trava para sempre. Por isso, nas duas threads, \textbf{primeiro se espera a disponibilidade, depois se pega a chave}.
\end{alerta}

\subsection{Vários produtores e vários consumidores}

O mesmo código funciona com várias threads produtoras e várias consumidoras, sem mudança nenhuma. \cmd{empty} e \cmd{full} continuam contando vagas e itens no total, e o \cmd{mutex} garante que duas produtoras nunca escrevam na mesma posição \cmd{hi} ao mesmo tempo. É por isso que o \cmd{mutex} é indispensável: sem ele, duas produtoras poderiam ler o mesmo \cmd{hi}, gravar na mesma posição e avançar o índice duas vezes, perdendo um item, que é a condição de corrida da Seção 3 de novo.

% ============================================================
\section{Deadlock, inanição e livelock}
% ============================================================

\subsection{As quatro condições do deadlock}

Em 1971, Coffman e colegas mostraram que um deadlock só acontece se \textbf{quatro} condições valem ao mesmo tempo:

\begin{center}
\begin{tabular}{@{}lR{0.4\linewidth}R{0.33\linewidth}@{}}
\toprule
\textbf{Condição} & \textbf{O que significa} & \textbf{No consumidor errado} \\
\midrule
Exclusão mútua & O recurso só pode ser usado por um de cada vez. & Só uma thread tem o \cmd{mutex}. \\
Posse e espera & Quem já tem um recurso espera por outro. & O consumidor tem o \cmd{mutex} e espera \cmd{full}. \\
Não preempção & Ninguém pode tomar um recurso à força. & O SO não tira o \cmd{mutex} do consumidor. \\
Espera circular & Existe um ciclo: A espera B, que espera A. & O consumidor espera o produtor, que espera o consumidor. \\
\bottomrule
\end{tabular}
\end{center}

Basta quebrar \textbf{uma} delas para o deadlock ser impossível. A ordem correta do produtor-consumidor quebra a \textbf{posse e espera}: nenhuma thread espera por \cmd{empty} ou \cmd{full} segurando o \cmd{mutex}.

\subsection{Como evitar na prática}

\begin{itemize}
    \item \textbf{Não dormir segurando uma trava} à espera de algo que outra thread só pode fazer com essa trava.
    \item \textbf{Ordem global de aquisição}: se um programa usa as travas A e B, todas as threads pegam sempre A antes de B. Assim não se forma o ciclo (quebra a espera circular).
    \item \textbf{Segurar travas pelo menor tempo possível}: fazer o trabalho demorado (cálculo, entrada e saída) fora da região crítica, como o comentário ``processamento do item fora do lock'' do \cmd{prod\_cons.c}.
\end{itemize}

\subsection{Inanição e livelock}

\begin{definicao}
\textbf{Inanição} (\emph{starvation}): uma thread fica esperando indefinidamente, não porque o sistema travou, mas porque as outras sempre passam na frente. O sistema progride; aquela thread, não.

\textbf{Livelock}: as threads não estão bloqueadas, continuam executando, mas reagem umas às outras de um jeito que ninguém avança. A imagem clássica são duas pessoas num corredor que dão um passo para o mesmo lado ao mesmo tempo, repetidamente.
\end{definicao}

A condição 4 de Tanenbaum (ninguém espera para sempre) é justamente a ausência de inanição. Peterson a garante para duas threads; um mutex comum, em geral, não promete a ordem de chegada.

% ============================================================
\section{Monitores em Java}
% ============================================================

\subsection{Por que monitores}

Semáforos resolvem o problema, mas deixam toda a disciplina nas mãos do programador. É fácil:

\begin{itemize}
    \item \textbf{esquecer} um \cmd{sem\_post}, e alguma thread dorme para sempre;
    \item \textbf{trocar a ordem} de dois \cmd{sem\_wait}, e o programa entra em deadlock;
    \item \textbf{proteger a coisa errada}, e a corrida continua.
\end{itemize}

\begin{definicao}[Monitor]
Um \textbf{monitor} é um objeto que junta os \textbf{dados compartilhados}, as \textbf{operações} que os acessam e a \textbf{exclusão mútua} dessas operações. Quem usa o objeto não precisa lembrar de pegar e devolver a chave: a linguagem faz isso.
\end{definicao}

Uma analogia comum: um banheiro com uma \textbf{chave única}. Quem tem a chave entra; quem chega depois espera na porta. Se quem está dentro descobre que falta papel, pode sair, \textbf{devolver a chave} e esperar ser avisado de que o papel chegou.

\subsection{Origem: Hoare e Brinch Hansen}

O conceito de monitor foi proposto por Per Brinch Hansen e por C. A. R. Hoare, no início dos anos 1970 (o artigo mais citado é o de Hoare, de 1974). A ideia tem duas partes:

\begin{enumerate}
    \item \textbf{Exclusão mútua automática}: só uma thread por vez executa dentro do monitor. O compilador, e não o programador, insere o ``pegar a chave'' e o ``devolver a chave''.
    \item \textbf{Variáveis de condição}: filas onde uma thread espera até que uma condição sobre os dados do monitor se torne verdadeira. As operações são \emph{wait} (esperar na condição, soltando o monitor) e \emph{signal} (acordar uma thread que espera na condição).
\end{enumerate}

Uma variável de condição \textbf{não guarda memória}, ao contrário do semáforo: um \emph{signal} sem ninguém esperando simplesmente se perde. Por isso a thread sempre testa a condição, sobre os dados do monitor, antes de esperar. Não há despertar perdido porque teste e espera acontecem dentro do monitor, com a chave na mão.

Em Java, todo objeto pode ser um monitor, com \textbf{uma única} variável de condição implícita: \cmd{wait}, \cmd{notify} e \cmd{notifyAll} operam nela.

\subsection{As ferramentas do Java}

\begin{center}
\begin{tabular}{@{}lR{0.68\linewidth}@{}}
\toprule
\textbf{Recurso} & \textbf{O que faz} \\
\midrule
\cmd{synchronized} & Para executar um método \cmd{synchronized}, a thread precisa da chave (lock) do objeto. Só uma thread por vez fica dentro dos métodos \cmd{synchronized} de um mesmo objeto. \\
\cmd{wait()} & Solta a chave do objeto e dorme até ser avisada. Só pode ser chamado por quem tem a chave (senão, \cmd{IllegalMonitorStateException}). \\
\cmd{notify()} & Acorda \textbf{uma} thread que dormia em \cmd{wait()} nesse objeto. Ela não recebe a chave na hora: precisa disputá-la de novo. \\
\cmd{notifyAll()} & Acorda \textbf{todas}; cada uma, ao conseguir a chave, confere a sua condição. \\
\cmd{Thread.sleep()} & Dorme por um tempo, mas \textbf{não solta a chave}. \\
\bottomrule
\end{tabular}
\end{center}

\subsection{Mailbox 3b: sem sincronização}

O exemplo das pastas \cmd{Monitor/1} e \cmd{Monitor/3} tem uma caixa de correio compartilhada por dois produtores e um consumidor. Dave envia ``Hello, world.'' e Bill envia ``Hot dog!''. Na versão 3b não há nenhuma sincronização:

\lstinputlisting[style=arquivojava, firstline=13, lastline=52, firstnumber=13]{\raiz/Monitor-20261002T192224Z-1-001/Monitor/1/Mailbox3b.java}

\cmd{storeMessage} escreve a mensagem \textbf{letra por letra} no mesmo array, com uma pequena pausa entre as letras. Se os dois produtores chamam ao mesmo tempo, as letras se misturam. Saída real:

\begin{lstlisting}[style=terminal]
My name is, Bill. I say, Hello, worl
My name is, Bill. I say, Hot dog!    .
My name is, Dave. I say, Het dog!
\end{lstlisting}

Bill ``disse'' a frase de Dave, a mensagem foi cortada, e ``Het dog!'' é metade de cada um. É a condição de corrida de novo, agora em Java. O \cmd{Thread.sleep} entre as letras só aumenta a janela em que a outra thread pode entrar: por isso aumentar \cmd{MAXPROCESSTIME} deixa o problema mais visível.

\subsection{Mailbox 3s: com monitor}

\lstinputlisting[style=arquivojava, firstline=6, lastline=49, firstnumber=6]{\raiz/Monitor-20261002T192224Z-1-001/Monitor/3/Mailbox3s.java}

\begin{itemize}
    \item Os dois métodos são \cmd{synchronized}: enquanto um produtor escreve, o outro produtor e o consumidor esperam pela chave. As mensagens saem \textbf{íntegras}.
    \item O consumidor, se não há mensagem, chama \cmd{wait()}: solta a chave (para que um produtor possa entrar) e dorme. Não fica perguntando ``chegou?'' o tempo todo.
    \item O produtor, ao terminar, chama \cmd{notify()} e acorda o consumidor.
\end{itemize}

\subsection{Por que while e não if}

O consumidor usa \cmd{while (youHaveMail == false) wait();}, e não um \cmd{if}. Motivos:

\begin{enumerate}
    \item Quem acorda \textbf{não recebe a chave na hora}: precisa disputá-la. Enquanto isso, outra thread pode entrar e mudar a condição.
    \item O Java permite \textbf{despertares espúrios}: uma thread pode sair do \cmd{wait()} sem que ninguém tenha chamado \cmd{notify}.
\end{enumerate}

Com \cmd{while}, a thread confere a condição de novo ao acordar e só segue quando ela é verdadeira de fato.

\subsection{Semântica de Hoare e semântica de Mesa}

Quando uma thread faz \emph{signal} e acorda outra, quem continua executando dentro do monitor?

\begin{itemize}
    \item \textbf{Hoare} (\emph{signal and wait}): a thread acordada entra \textbf{imediatamente} no monitor, e quem sinalizou espera. A condição que motivou o aviso ainda é verdadeira quando a acordada roda, e um \cmd{if} bastaria.
    \item \textbf{Mesa} (\emph{signal and continue}), do sistema Mesa da Xerox, em 1980: quem sinalizou \textbf{continua}, e a acordada só volta para a fila da chave. Até ela conseguir entrar, outra thread pode ter mudado tudo.
\end{itemize}

Java, pthreads (\cmd{pthread\_cond\_wait}) e quase todas as linguagens atuais usam a semântica de \textbf{Mesa}, que é mais simples de implementar e combina bem com despertares espúrios. Consequência prática, que vale para toda a vida: \textbf{espere sempre dentro de um \cmd{while}}.

\subsection{O limite da mailbox 3s}

\begin{alerta}[A 3s ainda pode perder mensagens]
\cmd{storeMessage} não espera a caixa esvaziar. Se dois produtores guardam mensagens antes de o consumidor ler, a primeira é \textbf{sobrescrita}. A versão 3s garante que cada mensagem lida está \textbf{íntegra}, mas não garante que toda mensagem enviada seja \textbf{entregue}.
\end{alerta}

Uma caixa de uma posição sem perdas é exatamente o produtor-consumidor com \cmd{N = 1}. Com monitor, ficaria assim:

\begin{lstlisting}[style=terminal]
guardar(mensagem), synchronized:       retirar(), synchronized:
    enquanto cheia: wait()                 enquanto vazia: wait()
    copia a mensagem                       copia a mensagem para o resultado
    cheia = verdadeiro                     cheia = falso
    notifyAll()                            notifyAll()
                                           devolve o resultado
\end{lstlisting}

Usa-se \cmd{notifyAll} porque, com dois produtores, um \cmd{notify} poderia acordar justamente o outro produtor (que veria a caixa cheia e voltaria a dormir) em vez do consumidor.

\subsection{Uma caixa de correio correta, em Java}

O pseudocódigo acima, escrito em Java e testado, está em \cmd{examples/dia07/caixa/}:

\lstinputlisting[style=arquivojava]{\raiz/examples/dia07/caixa/CaixaDeCorreio.java}

O \cmd{TesteCaixa.java} da mesma pasta põe dois produtores para enviar 20000 mensagens cada e confere o que chegou:

\begin{lstlisting}[style=terminal]
$ make run-caixa
Recebidas de Dave: 20000 | de Bill: 20000 (esperado: 20000 de cada)
\end{lstlisting}

Nenhuma mensagem se perde, e nenhuma sai misturada.

\subsection{Ferramentas prontas do Java}

Na prática, raramente se escreve \cmd{wait} e \cmd{notify} à mão. O pacote \cmd{java.util.concurrent} já oferece as estruturas desta aula, testadas e eficientes:

\begin{center}
\begin{tabular}{@{}lR{0.6\linewidth}@{}}
\toprule
\textbf{Classe} & \textbf{Corresponde a} \\
\midrule
\cmd{Semaphore} & Os semáforos do C (\cmd{acquire} = \cmd{sem\_wait}, \cmd{release} = \cmd{sem\_post}). É o que a simulação Produção de dinheiro usa. \\
\cmd{ReentrantLock} e \cmd{Condition} & Um monitor explícito, com quantas variáveis de condição se quiser (por exemplo, uma para ``não cheio'' e outra para ``não vazio''). \\
\cmd{ArrayBlockingQueue} & O buffer limitado inteiro: \cmd{put} espera vaga e \cmd{take} espera item. É o produtor-consumidor pronto. \\
\cmd{AtomicInteger} & O contador atômico: \cmd{incrementAndGet} é o \cmd{atomic\_fetch\_add} do Java. \\
\bottomrule
\end{tabular}
\end{center}

\subsection{Semáforos ou monitores?}

\begin{center}
\begin{tabular}{@{}R{0.2\linewidth}R{0.36\linewidth}R{0.36\linewidth}@{}}
\toprule
& \textbf{Semáforo} & \textbf{Monitor} \\
\midrule
Exclusão mútua & Feita à mão, com um semáforo em 1 & Automática \\
Esperar condição & Contadores (\cmd{empty}, \cmd{full}) & \cmd{wait} dentro de \cmd{while} \\
Guarda avisos? & Sim: o valor acumula os \cmd{post} & Não: \cmd{notify} sem ninguém esperando se perde \\
Onde fica a disciplina & Espalhada pelo código que usa & Concentrada no objeto \\
Erros típicos & Esquecer um \cmd{post}, trocar a ordem & Usar \cmd{if}, esquecer \cmd{notifyAll} \\
\bottomrule
\end{tabular}
\end{center}

Os dois têm o mesmo poder: dá para implementar um com o outro. A escolha é de organização: o monitor deixa os erros mais difíceis de cometer, porque a chave é pega e devolvida pela linguagem.

% ============================================================
\section{Demonstração visual: Produção de dinheiro}
% ============================================================

\subsection{A simulação}

A simulação ``Produção de dinheiro'' mostra o produtor-consumidor acontecendo na tela. O \textbf{banco} é o produtor (cria notas) e o \textbf{carro-forte} é o consumidor (retira notas). Por baixo, ela usa threads Java e três semáforos, como o \cmd{prod\_cons.c}.

\begin{lstlisting}[style=terminal]
$ make run-simulation
\end{lstlisting}

A janela precisa de JDK 21 ou mais novo e de Maven; na primeira execução, o Maven baixa o JavaFX pela internet.

\begin{center}
\begin{tabular}{@{}ll@{}}
\toprule
\textbf{Na tela} & \textbf{No \cmd{prod\_cons.c}} \\
\midrule
10 posições & \cmd{N} \\
W (próxima escrita) & \cmd{hi} \\
R (próxima leitura) & \cmd{lo} \\
itens disponíveis & \cmd{full} \\
vagas livres & \cmd{empty} \\
``Produtor: Aguardando vaga'' & dormindo em \cmd{sem\_wait(\&empty)} \\
``Consumidor: Aguardando item'' & dormindo em \cmd{sem\_wait(\&full)} \\
\bottomrule
\end{tabular}
\end{center}

\subsection{Quatro experimentos}

Para cada um, faça a previsão antes de olhar.

\subsubsection{Experimento 1: ritmos iguais}
Produtor e consumidor em 2 itens/s. O buffer fica oscilando com poucas notas, e W e R andam quase juntos. Há quase sempre vaga e quase sempre item: ninguém espera muito.

\subsubsection{Experimento 2: produtor rápido}
Produtor em 10, consumidor em 0,2. O buffer enche e aparece ``Produtor: Aguardando vaga: buffer cheio''. No C, o produtor está dormindo em \cmd{wait\_sem(\&empty)} com \cmd{empty = 0}. Ele \textbf{não} segura o \cmd{mutex}, e por isso o consumidor continua conseguindo retirar.

\subsubsection{Experimento 3: consumidor rápido}
Produtor em 0,2, consumidor em 10. O buffer esvazia e aparece ``Consumidor: Aguardando item: buffer vazio''. O consumidor dorme em \cmd{wait\_sem(\&full)} com \cmd{full = 0}; cada nota nova faz \cmd{sem\_post(\&full)} e o acorda.

\subsubsection{Experimento 4: a volta do buffer e a pausa}
Com ritmos intermediários, W chega à posição 9 e volta para a 0: é o \cmd{(hi + 1) \% N}. Pausando, itens disponíveis mais vagas livres dá sempre 10: é o \cmd{empty + full = N}.

\begin{exemplo}[O que a simulação mostra]
O \cmd{mutex} garante que ninguém atrapalha. \cmd{empty} e \cmd{full} garantem que ninguém trabalha sem ter o que precisa.
\end{exemplo}

% ============================================================
\section{Atividade prática}
% ============================================================

\subsection{O que fazer}

Abra \cmd{examples/dia07/prod\_cons\_atividade.c}. É o \cmd{prod\_cons.c} com seis lacunas:

\begin{lstlisting}[style=codigo]
sem_init(&mutex, 0, 1);
sem_init(&empty, 0, @/* 1: vagas no inicio */@);
sem_init(&full,  0, @/* 2: itens no inicio */@);

/* produtor */                   /* consumidor */
wait_sem(@/* 3 */@);                 wait_sem(@/* 5 */@);
wait_sem(&mutex);                wait_sem(&mutex);
insert_item(item);               item = remove_item();
sem_post(&mutex);                sem_post(&mutex);
sem_post(@/* 4 */@);                 sem_post(@/* 6 */@);
\end{lstlisting}

O \cmd{mutex} já está pronto e a ordem das linhas não deve mudar. As perguntas que guiam:

\begin{itemize}
    \item Quantas vagas e quantos itens existem quando o programa começa?
    \item Antes de mexer no buffer, o que cada thread precisa ter?
    \item Depois de mexer, quem precisa ficar sabendo, e de quê?
\end{itemize}

\subsection{Compilar e testar}

\begin{lstlisting}[style=terminal]
$ gcc -std=c11 -Wall -Wextra -Wpedantic -Werror \
      examples/dia07/prod_cons_atividade.c -o prod-cons-aluno -pthread
$ ./prod-cons-aluno
...
Fim: 40 itens produzidos e consumidos; buffer vazio.
\end{lstlisting}

Antes de preencher, o esqueleto não compila. Isso é proposital.

\subsection{Se der errado}

\begin{center}
\begin{tabular}{@{}R{0.42\linewidth}R{0.5\linewidth}@{}}
\toprule
\textbf{Sintoma} & \textbf{O que significa} \\
\midrule
O programa para de imprimir e não termina (saia com \cmd{Ctrl + C}) & Alguma thread dorme esperando um semáforo que ninguém vai liberar. \\
\cmd{Assertion `count < N' failed} & O produtor inseriu sem haver vaga: \cmd{empty} não está diminuindo de verdade. \\
\cmd{Assertion `count > 0' failed} & O consumidor retirou sem haver item: \cmd{full} está maior do que deveria. \\
\bottomrule
\end{tabular}
\end{center}

\subsection{Perguntas para responder em dupla}

\begin{enumerate}
    \item Por que o produtor espera \cmd{empty} \textbf{antes} de pegar o \cmd{mutex}? O que poderia acontecer se pegasse o \cmd{mutex} primeiro?
    \item Se o seu programa travou ou abortou em alguma tentativa, qual lacuna estava errada e por que isso causou aquele sintoma?
\end{enumerate}

% ============================================================
\section{Resumo}
% ============================================================

\begin{center}
\begin{tabular}{@{}lR{0.66\linewidth}@{}}
\toprule
\textbf{Conceito} & \textbf{Em uma frase} \\
\midrule
Thread & Fluxo de execução que divide a memória com as outras threads do processo. \\
Condição de corrida & O resultado depende da ordem em que as threads foram intercaladas. \\
Região crítica & Trecho que acessa o recurso compartilhado. \\
Exclusão mútua & No máximo uma thread por vez na região crítica. \\
Peterson & Exclusão mútua para duas threads, com espera ocupada. \\
Semáforo & Contador com \cmd{wait} e \cmd{post}; quem não pode seguir dorme. \\
Produtor-consumidor & \cmd{mutex} para exclusão mútua; \cmd{empty} e \cmd{full} para disponibilidade. \\
Instrução atômica & Lê e escreve a memória numa operação indivisível, mesmo com vários núcleos. \\
Spinlock & Trava em que quem espera fica girando; boa para esperas curtíssimas. \\
Deadlock & Cada thread espera algo que só a outra pode dar; exige as quatro condições de Coffman. \\
Inanição & Uma thread espera indefinidamente enquanto as outras passam na frente. \\
Variável de condição & Fila de espera dentro de um monitor; não guarda avisos. \\
Monitor & Dados, operações e exclusão mútua no mesmo objeto; \cmd{wait} e \cmd{notify} para esperar condições. \\
\bottomrule
\end{tabular}
\end{center}

\begin{definicao}[As duas aulas em uma frase]
Processos isolam a memória e precisam de pipe para conversar. Threads compartilham a memória e precisam de sincronização para não se atrapalhar.
\end{definicao}

% ============================================================
\section{Exercícios de fixação}
% ============================================================

\begin{enumerate}
    \item Por que o resultado de \cmd{corrida.c} muda a cada execução?
    \item Escreva uma intercalação dos três passos de \cmd{contador++} em duas threads que dê o resultado \textbf{correto} (2) e outra que dê o resultado errado (1).
    \item No produtor-consumidor com \cmd{N = 5}, depois de três produções e um consumo completos, quanto valem \cmd{empty}, \cmd{full} e \cmd{count}?
    \item Por que Peterson não obriga as threads a se alternarem (uma vez cada)?
    \item E se o semáforo \cmd{mutex} fosse iniciado com 0 em vez de 1?
    \item No produtor, trocar a ordem dos dois \cmd{sem\_post} (avisar \cmd{full} antes de devolver o \cmd{mutex}) causa deadlock?
    \item Por que \cmd{wait()} só pode ser chamado dentro de um método \cmd{synchronized}?
    \item Na mailbox completa com dois produtores e um consumidor, por que usar \cmd{notifyAll()} e não \cmd{notify()}?
    \item Se o consumidor da mailbox 3s usasse \cmd{Thread.sleep(100)} em vez de \cmd{wait()} dentro do \cmd{while}, o que aconteceria?
    \item Das 20 intercalações possíveis de dois \cmd{contador++}, quais dão o resultado certo? Por quê?
    \item Com \cmd{-O2}, o \cmd{contador = contador + 1} vira uma única instrução \cmd{addl}. Por que a corrida continua?
    \item Qual das quatro condições de uma boa solução a alternância estrita viola? Dê um exemplo.
    \item No deadlock do consumidor que pega o \cmd{mutex} antes de esperar \cmd{full}, identifique as quatro condições de Coffman.
    \item Num monitor com semântica de Hoare, o \cmd{while} antes de \cmd{wait} seria necessário? E em Java?
    \item Por que um semáforo não sofre o problema do despertar perdido, mas o par \cmd{dormir()}/\cmd{acordar()} ingênuo sofre?
\end{enumerate}

\subsection*{Respostas}

\begin{enumerate}
    \item Porque a ordem em que as threads são intercaladas depende do escalonador, da carga da máquina e de quantos núcleos estão livres, e isso muda a cada execução. Quantos incrementos se perdem depende de quantas vezes as leituras das duas threads caíram ``no meio'' uma da outra.
    \item Correta: A lê 0, A calcula 1, A escreve 1, B lê 1, B calcula 2, B escreve 2. Errada: a da Seção~\ref{sec:intercalacao}, em que as duas leem 0 antes de alguma escrever.
    \item \cmd{empty = 3}, \cmd{full = 2}, \cmd{count = 2}. Três produções levam \cmd{empty} de 5 a 2 e \cmd{full} de 0 a 3; o consumo devolve uma vaga e retira um item.
    \item Porque a espera só acontece se o \textbf{outro} também estiver interessado. Se a outra thread não quer entrar, \cmd{interessado[outro]} é falso e a thread entra quantas vezes quiser seguidas.
    \item Ninguém conseguiria pegar a chave: o primeiro \cmd{sem\_wait(\&mutex)} de qualquer thread dormiria para sempre, e o programa travaria logo no começo.
    \item Não. O consumidor pode acordar um pouco antes e ter de esperar o \cmd{mutex} até o produtor devolvê-lo, mas o produtor devolve logo em seguida, sem esperar nada. Funciona; só muda um pouco o desempenho. O deadlock vem de \textbf{dormir} segurando a chave, não da ordem dos \cmd{post}.
    \item Porque \cmd{wait()} precisa \textbf{soltar} a chave, e só quem tem a chave pode soltá-la. Além disso, testar a condição e começar a esperar precisa acontecer sem que outra thread mude a condição no meio. Fora de \cmd{synchronized}, o Java lança \cmd{IllegalMonitorStateException}.
    \item Com \cmd{notify()}, o produtor que acabou de guardar pode acordar o \textbf{outro produtor}, que vê a caixa cheia e volta a dormir. O consumidor continua dormindo, e todos podem acabar esperando para sempre. Com \cmd{notifyAll()}, todos acordam e cada um confere a sua condição: o consumidor certamente acorda.
    \item O consumidor dormiria \textbf{segurando a chave}, porque \cmd{sleep} não a solta. Nenhum produtor conseguiria entrar em \cmd{storeMessage} para colocar uma mensagem, e o consumidor ficaria acordando e dormindo para sempre com a caixa vazia.
    \item Só AAABBB e BBBAAA, em que uma thread faz os três passos antes de a outra começar. Em qualquer outra, as duas leem o valor antes de alguma escrever, e uma escrita apaga a outra.
    \item Porque uma instrução não é atômica entre núcleos: por dentro, o processador lê, soma e escreve, e outro núcleo pode ler no meio. Só o prefixo \cmd{lock} (gerado por \cmd{atomic\_fetch\_add}) torna a operação indivisível.
    \item A condição 3: uma thread fora da região crítica não pode impedir outra de entrar. Se a thread 0 quer entrar duas vezes seguidas, precisa esperar a thread 1 entrar e sair, mesmo que a thread 1 nem queira entrar.
    \item Exclusão mútua: só um tem o \cmd{mutex}. Posse e espera: o consumidor tem o \cmd{mutex} e espera \cmd{full}. Não preempção: ninguém tira o \cmd{mutex} dele. Espera circular: o consumidor espera um item do produtor, e o produtor espera o \cmd{mutex} do consumidor.
    \item Com Hoare, a thread acordada entra imediatamente, com a condição garantida, e um \cmd{if} bastaria (desde que não houvesse despertares espúrios). Java usa a semântica de Mesa e permite despertares espúrios: o \cmd{while} é obrigatório.
    \item Porque o semáforo guarda os avisos no seu valor: um \cmd{sem\_post} feito antes do \cmd{sem\_wait} não se perde, apenas soma 1, e o \cmd{sem\_wait} seguinte passa direto. Além disso, o teste do valor e o bloqueio acontecem numa operação indivisível. No par ingênuo, o teste e o \cmd{dormir()} são separados, e o \cmd{acordar()} feito no meio não deixa registro.
\end{enumerate}

% ============================================================
\section*{Para saber mais}
\addcontentsline{toc}{section}{Para saber mais}
% ============================================================

\begin{itemize}
    \item Páginas de manual: \cmd{man 7 pthreads}, \cmd{man 3 pthread\_create}, \cmd{man 3 sem\_wait}, \cmd{man 3 sem\_init}. Também em \url{https://man7.org/linux/man-pages/}.
    \item Contrato de \cmd{wait}, \cmd{notify} e despertares espúrios: documentação de \cmd{java.lang.Object}, em \url{https://docs.oracle.com/en/java/javase/21/docs/api/java.base/java/lang/Object.html}.
    \item TANENBAUM, A. S.; BOS, H. \emph{Sistemas Operacionais Modernos}. 4. ed. São Paulo: Pearson, 2016. Seção 2.3: comunicação entre processos (condições de corrida, regiões críticas, Peterson, semáforos, monitores).
    \item SILBERSCHATZ, A.; GALVIN, P. B.; GAGNE, G. \emph{Fundamentos de Sistemas Operacionais}. 9. ed. Rio de Janeiro: LTC, 2015. Capítulo 5: sincronização de processos.
\end{itemize}

Em outras edições, o assunto continua nos capítulos sobre sincronização, mas a numeração pode mudar.

\begin{itemize}
    \item DOWNEY, A. B. \emph{The Little Book of Semaphores}. Green Tea Press, livre, em \url{https://greenteapress.com/wp/semaphores/}. Dezenas de problemas de sincronização resolvidos com semáforos, em inglês.
    \item HERLIHY, M.; SHAVIT, N. \emph{The Art of Multiprocessor Programming}. Morgan Kaufmann. Capítulo 2: exclusão mútua, com a prova completa de Peterson.
    \item GOETZ, B. et al. \emph{Java Concurrency in Practice}. Addison-Wesley, 2006. Monitores, o pacote \cmd{java.util.\allowbreak concurrent} e os erros mais comuns em Java.
\end{itemize}

\end{document}
